




\usepackage{amsfonts,latexsym} % para tener disponibilidad de diversos simbolos


\usepackage{array,colortbl}
%\usepackage{ifpdf}


%\usepackage{rotating}
\usepackage{cite}
\usepackage{stfloats}
\usepackage{listings}
\usepackage{fancyhdr}


%\usepackage{wrapfig}
\usepackage{array}
\usepackage[symbol]{footmisc}
%\usepackage{multirow}
%\usepackage{adjustbox}
%\usepackage{nccmath}
% BOXES 
\usepackage{tikz}
\usepackage[most]{tcolorbox}
\tcbset{blue box/.style={
    enhanced, fonttitle=\bfseries,
    colback=white, colframe=accdarklight,
    coltitle=white, colbacktitle=accdark,
    attach boxed title to top left={xshift=0.3cm,
                                    yshift*=-\tcboxedtitleheight/2},
    boxed title style={
      before upper=\hspace*{0.5cm}, % reserve space for the image
      overlay={
       \node at ([xshift=0.5cm]frame.west) {};
        
      }
    }
  }
}
\newtcolorbox{mybox}[1][]{blue box, #1}





% COLORS

\usepackage{xcolor}
\definecolor{accdark}{HTML}{102542}
\definecolor{accdarklight}{HTML}{204983}
\definecolor{acclight}{HTML}{FFCB77}
\definecolor{bck}{HTML}{fcfffc}
\definecolor{urgent}{HTML}{A02C3C}
\definecolor{colorthree}{HTML}{107E7D}
\definecolor{colorfour}{HTML}{105260}
\definecolor{myred}{HTML}{af5655}
\definecolor{mygreen}{HTML}{377771}
\definecolor{myblue}{HTML}{164254}
\definecolor{myyellow}{HTML}{D5BF86}
\definecolor{mygray}{HTML}{8AA1B1}
\definecolor{lightgray}{gray}{0.9}
\definecolor{bckdarker}{HTML}{dfeedf}
\definecolor{urgentlight}{HTML}{CF4F60}
\definecolor{colortwo}{HTML}{3AB09E}
\definecolor{codebg}{HTML}{1f1f1f}
\definecolor{boxTitle}{HTML}{B1CF5F}
\definecolor{rose}{HTML}{CC7E85}
\definecolor{rose-d}{HTML}{83343C}
\definecolor{rose-l}{HTML}{DCA7AC}
\definecolor{boxBackground}{HTML}{DEF4C6}
\definecolor{boxFrame}{HTML}{004346}
%%
\definecolor{grays}{HTML}{4F646F}
\definecolor{blush}{HTML}{EA638C}
\definecolor{reds}{HTML}{7D1D3F}
\definecolor{ash}{HTML}{AEB4A9}
\definecolor{blues}{HTML}{2C4251}
\definecolor{purp}{HTML}{850F78}
\definecolor{orang}{HTML}{F59121}
\definecolor{yello}{HTML}{fae99e}

% Glossary 

\newglossary*{dinge}{Was ist was?}
\glsaddkey
 {formB}
 {}
 {\glsentryB}
 {\GlsentryB}
 {\glsB}
 {\GlsB}
 {\GLSB}
 \glsaddkey
 {formF}
 {}
 {\glsentryF}
 {\GlsentryF}
 {\glsF}
 {\GlsF}
 {\GLSF}
 
\newglossary[slg]{symbolslist}{syi}{syg}{Symbole und Begriffe} % create add. symbolslist

\glsaddkey{unit}{\glsentrytext{\glslabel}}{\glsentryunit}{\GLsentryunit}{\glsunit}{\Glsunit}{\GLSunit}
\glsaddkey{verb}{\glsentrytext{\glslabel}}{\glsentryverb}{\GLsentryverb}{\glsverb}{\Glsverb}{\GLSverb}





\newcommand{\charfunc}[1]{\glsentrytext{charfunc} mit $A = #1$, also $\chi_{#1}$}



\newcommand{\mapsfrom}{\mathrel{\style{display:inline-block; transform:scale(-1,1);}{\mapsto}}}
\newcommand{\X}{$X$}
\usepackage{amsthm}
%\theoremstyle{definition}
\newtheoremstyle{mytheoremstyle} % name
    {\topsep}                    % Space above
    {\topsep}                    % Space below
    {}                   % Body font
    {}                           % Indent amount
    {\bfseries}                   % Theorem head font
    {}                          % Punctuation after theorem head
    {1.5em}                       % Space after theorem head
    {\thmname{#1}\thmnumber{ #2}\thmnote{ #3}}  % Theorem head spec (can be left empty, meaning ‘normal’)


\theoremstyle{mytheoremstyle}
\newtheorem{test}{Test}[section]
\newtheorem{bsp}[test]{Beispiel}
\newtheorem{info}[test]{Zusatzinformation}
\newtheorem{reg}[test]{Regel}
\newtheorem{kor}[test]{Korollar}
\newtheorem{satz}[test]{Satz}
\tcolorboxenvironment{satz}{
  colback=white,
  colframe=ash,
  boxrule=0.8pt,
  arc=2mm,
  top=1mm,
  bottom=1mm,
  left=1mm,
  right=1mm,
  enhanced,
  breakable,
  parbox=false
}
\newtheorem{lem}[test]{Lemma}
\tcolorboxenvironment{lem}{
  colback=white,
  colframe=ash,
  boxrule=0.8pt,
  arc=2mm,
  top=1mm,
  bottom=1mm,
  left=1mm,
  right=1mm,
  enhanced,
  breakable,
  parbox=false
}
\newtheorem{defi}[test]{Definition}
\tcolorboxenvironment{defi}{
  colback=white,
  colframe=ash,
  boxrule=0.8pt,
  arc=2mm,
  top=1mm,
  bottom=1mm,
  left=1mm,
  right=1mm,
  enhanced,
  breakable,
  parbox=false
}
\newtheorem{fr}[test]{Frage}
\newtheorem{beo}[test]{Beobachtung}
\tcolorboxenvironment{bsp}{
  colback=white,
  colframe=ash,
  boxrule=0.8pt,
  arc=2mm,
  top=1mm,
  bottom=1mm,
  left=1mm,
  right=1mm,
  enhanced,
  breakable,
  parbox=false
}
\newtheorem{prinz}[test]{Prinzip}
\newtheorem{bem}[test]{Bemerkung}
\tcolorboxenvironment{bem}{
  colback=gray!10,
  colframe=ash,
  boxrule=0.8pt,
  arc=2mm,
  top=1mm,
  bottom=1mm,
  left=1mm,
  right=1mm,
  enhanced,
  breakable,
  parbox=false
}
\newtheorem{anmerkung}{Anmerkung}[]
\newtheorem{folg}[test]{Folgerung}
\newtheorem{erg}[test]{Ergänzung}
\tcolorboxenvironment{erg}{
  colback=yellow!5,
  colframe=ash,
  boxrule=0.8pt,
  arc=2mm,
  top=1mm,
  bottom=1mm,
  left=1mm,
  right=1mm,
  enhanced,
  breakable,
  parbox=false
}
\newtheorem{er}[test]{Erinnerung}
\tcolorboxenvironment{er}{
  colback=blue!10,
  colframe=ash,
  boxrule=0.8pt,
  arc=2mm,
  top=1mm,
  bottom=1mm,
  left=1mm,
  right=1mm,
  enhanced,
  breakable,
  parbox=false
}
\newtheorem{prob}[test]{Problem}
\newtheorem{eig}[test]{Eigenschaften}
\newcommand{\anm}[2]{
\color{Blue} \noindent
\begin{anmerkung}  \textit{zu "#1"}  \newline
#2
\end{anmerkung} \color{black}
}

\newcommand{\regel}[2]{
\begin{reg} \textbf{#1}  \newline
#2
\end{reg}
}
\newcommand{\prin}[2]{
\begin{prinz} \textbf{#1}  \newline
#2
\end{prinz}
}
\newcommand{\rsp}[3]{
\noindent \textbf{#1} \pageref{rsp:#2}  

\noindent #3 

}
\newcommand{\txc}[1]{
\textcolor{Blue}{#1}
}
\newcommand{\txo}[1]{
\textcolor{Orange}{#1}
}

\newcommand{\p}[0]{\mathbb{P}}
\newcommand{\za}[1]{\mathbb{Z} / #1 \mathbb{Z}}
\newcommand{\zat}[0]{\mathbb{Z} / m \mathbb{Z}}
\newcommand{\zatx}[1]{\mathbb{Z} / #1 \mathbb{Z}[x]}
\usepackage[onehalfspacing]{setspace} 
\renewcommand{\thefootnote}{\roman{footnote}}
\newcommand{\limnin}[0]{\lim_{n \rightarrow \infty}}
\newcommand{\intab}[0]{\int_a^b}
\newcommand{\dx}[0]{\,dx}
\newcommand{\sumEN}[0]{\sum_{i=1}^n}

\newcommand{\explainvkr}[0]{\footnote{$v$: Anzahl der Elemente in der Menge $X$} \footnote{$k$: Anzahl der Elemente pro Block} \footnote{$r$: in wie vielen Blöcken ein Element vorkommen soll}}
\newcommand{\explainvkrb}[0]{\explainvkr \footnote{$b$: Anzahl der Blöcke im Design}}
\newcommand{\explaindesign}[0]{\footnote{$\mathcal{B}, \mathcal{C}, ...$: damit sind Designs gemeint}}
\newcommand{\explainrs}[0]{\footnote{$r_s, r_t, ...$: t bzw. s ist die Größe der Teilmengen von X, die gleich oft in allen $r_t$ Blöcken eines $t$-Designs vorkommen soll}}

\usepackage{todonotes}
\newcommand{\bluenote}[1]{\todo[inline, inlinewidth=0.9\linewidth, bordercolor=Blue, backgroundcolor=white]{\textcolor{Blue}{#1}}}
\newcommand{\orangenote}[1]{\todo[inline, inlinewidth=0.9\linewidth, bordercolor=Orange, backgroundcolor=white]{\textcolor{Orange}{#1}} }
%\ifvmode\leavevmode\fi

\newcommand{\kl}[1]{\{ #1 \}}
\newcommand{\txt}[1]{\quad \text{#1}}
\newcommand{\faoio}[0]{ \quad \forall \glsunit{elementarereignis} \in \glsunit{eraum} }
\newcommand{\oio}[0]{\glsunit{elementarereignis} \in \glsunit{eraum}}
\newcommand{\wraum}[0]{\glsunit{wraum}}
\newcommand{\oub}[2]{\txo{
     \underbrace{
     \textcolor{black}{#1}
     }_{#2}
     }}
\newcommand{\bub}[2]{\txc{
\underbrace{
\textcolor{black}{#1}
}_{#2}
}}
\newcommand{\oob}[2]{\txo{
\overbrace{
\textcolor{black}{#1}
}^{#2}
}}
\newcommand{\ma}[1]{\(#1\)}
\newcommand{\mab}[1]{\[#1\]}
\newcommand{\E}[1]{\mathbb{E}[#1]}
\newcommand{\dichte}[0]{\Gls{dichte} \index{Dichte}}

\newglossarystyle{symbunitlong}{%
\setglossarystyle{long3col}% base this style on the list style
\renewenvironment{theglossary}{% Change the table type --> 3 columns
  \begin{longtable}{lp{1\glsdescwidth}>{\centering\arraybackslash}p{2cm}}}%
  {\end{longtable}}%
%
\renewcommand*{\glossaryheader}{%  Change the table header
  \bfseries Gegenstand & \bfseries Beschreibung & \bfseries Symbol \\
  \hline
  \endhead}
\renewcommand*{\glossentry}[2]{%  Change the displayed items
\glstarget{##1}{\glossentryname{##1}} %
& \glossentrydesc{##1}% Description
& \glsunit{##1}  \tabularnewline
}
}

\newcommand{\pe}[0]{\glsunit{wmaß}}
\newcommand{\bernoulli}[0]{\Gls{bernoulli} \index{Bernoulli-Experiment}}
\newcommand{\binomv}[0]{\Gls{binomv} \index{Binomialverteilung}}
\newcommand{\ele}[0]{\glsunit{elementarereignis}}
\newcommand{\eraum}[0]{\glsunit{eraum}}
\newcommand{\zv}[0]{\Gls{zv}}